% Unit 1 free-response set -- 2 long (10) + 3 short (4) = 32 pts.
%
% The CED track's FRQ sets deliberately CROSS chapters where the unit does.
% Unit 1 draws on Zumdahl Ch3 (moles, mass spectra), Ch7 (structure, PES,
% trends) and Ch2/Ch8 (valence, ionic compounds); the chapter-track sets treat
% those separately, so the questions here integrate them instead of repeating
% them. No item duplicates a chapter FRQ set.
\documentclass{apchem-exam}

\apcunit{1}
\apcunittitle{Atomic Structure and Properties}
\apcdoctitle{Free-Response Set}
\apcsubtitle{CED 1.1--1.8 \textbullet\ 5 questions \textbullet\ 32 points \textbullet\ 50 minutes}

\begin{document}
\apctitleblock

\begin{apcnote}[Scope --- two CED exclusions apply]
Covers \ced{1.1} through \ced{1.8}. Unit 1 is worth
\SI{7}{\percent}--\SI{9}{\percent} of the exam.

Topic \ced{1.5} excludes \textbf{assigning quantum numbers to individual
electrons}, and \ced{1.7} excludes \textbf{electron configurations of the
aufbau exceptions} (\ce{Cr}, \ce{Cu}). Neither is asked below.

Trend explanations must be \textbf{Coulombic} --- nuclear charge,
distance, shielding. The rubrics award zero for answers resting on
``stability'' or ``wants a full octet''.
\end{apcnote}

\examsection{II}{Free Response}{5 questions \textbullet\ 32 points \textbullet\ 50 minutes}{\calculatorok}

% =====================================================================
\frq{long}{10}{\ced{1.1} \ced{1.2} \textbullet\ mass spectrum to mole count}

Boron occurs as two isotopes. Its mass spectrum shows peaks at
\SI{10.013}{\atomicmassunit} (\SI{19.9}{\percent}) and
\SI{11.009}{\atomicmassunit} (\SI{80.1}{\percent}).

\begin{parts}
  \item Calculate the average atomic mass of boron. \workspace[2.6cm]
        \keyonly{\begin{apcanswer}$(10.013)(0.199) + (11.009)(0.801)
        = 1.993 + 8.818 = \mathbf{10.81}$~u, matching the periodic
        table.\end{apcanswer}}
  \item State which peak is taller and explain what peak height
        represents. \answerlines[2]{The \SI{11.009}{} peak. Peak height is
        proportional to the relative abundance of that isotope, and
        \ce{^{11}B} accounts for \SI{80.1}{\percent} of boron atoms.}
  \item Calculate the number of moles in a \SI{5.00}{\gram} sample of
        boron, and the number of atoms it contains. \workspace[2.8cm]
        \keyonly{\begin{apcanswer}$n = 5.00/10.81 = \mathbf{0.463}$~mol

        atoms $= 0.4625 \times 6.022\times10^{23} =
        \mathbf{2.79\times10^{23}}$\end{apcanswer}}
  \item Explain why the average atomic mass must be used for the mole
        calculation in (c) rather than either isotope mass.
        \answerlines[3]{A real sample contains both isotopes in their
        natural proportions, so its average particle mass is the
        abundance-weighted value. Using \num{10.013} or \num{11.009} would
        assume the sample were isotopically pure, which it is not.}
  \item A second element has two isotopes of nearly equal abundance and an
        average atomic mass almost exactly halfway between them. State
        what this implies about the relative peak heights.
        \answerlines[2]{The two peaks are of nearly equal height. The
        average sits midway only when the two isotopes contribute almost
        equally, which means almost equal abundances.}
\end{parts}

\begin{rubric}
  \rpoint{2}{(a) 1 pt weighting each mass by its fractional abundance;
    1 pt \num{10.81}. A plain mean of the two masses gives \num{10.51} and
    earns 0}
  \rpoint{2}{(b) 1 pt the \ce{^{11}B} peak; 1 pt height means relative
    abundance. Confusing height with mass earns 0}
  \rpoint{3}{(c) 1 pt \num{0.463}~mol; 1 pt multiplying by Avogadro's
    number; 1 pt \SI{2.79e23}{}}
  \rpoint{2}{(d) 1 pt the sample contains both isotopes; 1 pt the average
    reflects their proportions}
  \rpoint{1}{(e) nearly equal peak heights}
\end{rubric}

% =====================================================================
\frq{long}{10}{\ced{1.5} \ced{1.6} \ced{1.7} \ced{1.8} \textbullet\ PES to configuration to trend}

A photoelectron spectrum of an unknown element from period 3 shows four
peaks with relative heights $2$, $2$, $6$, $2$ and a fifth of height $1$,
listed from highest binding energy to lowest.

\begin{parts}
  \item Identify the element and write its complete electron
        configuration. \workspace[2.6cm]
        \keyonly{\begin{apcanswer}The heights give the electron count in
        each subshell: $2+2+6+2+1 = 13$ electrons, so $Z = 13$ and the
        element is \textbf{aluminium}.

        \ce{1s^2 2s^2 2p^6 3s^2 3p^1}\end{apcanswer}}
  \item State which peak lies at the highest binding energy and justify it
        in Coulombic terms. \answerlines[3]{The height-2 \ce{1s} peak.
        Those electrons sit closest to the nucleus and are shielded by no
        other electrons, so they experience the largest effective nuclear
        charge at the smallest distance; the attraction is strongest and
        the most energy is needed to remove one.}
  \item The first ionization energies of magnesium and aluminium are 738
        and 578 \si{\kilo\joule\per\mole} respectively. Magnesium's is
        \emph{higher}, even though aluminium has the greater nuclear
        charge. Explain. \workspace[3.2cm]
        \keyonly{\begin{apcanswer}The electron removed comes from a
        \emph{different subshell} in each case.

        Magnesium is \ce{3s^2}: its outermost electron is a \ce{3s}
        electron. Aluminium is \ce{3s^2 3p^1}: its outermost electron is
        in a \ce{3p} orbital, which lies at higher energy and, on average,
        further from the nucleus. The \ce{3p} electron is also shielded by
        the filled \ce{3s} sublevel beneath it.

        Greater distance plus greater shielding means a smaller net
        attraction, so the aluminium electron is easier to remove despite
        the extra proton.\end{apcanswer}}
  \item State the ion aluminium forms and give the formula of the compound
        it forms with oxygen. \answerlines[2]{\ce{Al^3+}; with \ce{O^2-}
        the neutral formula is \textbf{\ce{Al2O3}}.}
\end{parts}

\begin{rubric}
  \rpoint{3}{(a) 1 pt summing the peak heights to 13 electrons; 1 pt
    aluminium; 1 pt the full configuration}
  \rpoint{3}{(b) 1 pt the \ce{1s} peak; 1 pt closest to the nucleus /
    unshielded; 1 pt greatest effective nuclear charge. ``It is the
    innermost'' alone earns 1}
  \rpoint{2}{(c) 1 pt identifying the \ce{3s}-versus-\ce{3p} difference;
    1 pt greater distance and/or shielding of the \ce{3p} electron.
    An answer arguing only from nuclear charge predicts the wrong order
    and earns \textbf{0}}
  \rpoint{2}{(d) 1 pt \ce{Al^3+}; 1 pt \ce{Al2O3}}
\end{rubric}

% =====================================================================
\frq{short}{4}{\ced{1.3} \textbullet\ combustion analysis}

A \SI{0.500}{\gram} sample of a compound containing only carbon, hydrogen
and oxygen is burned completely, producing \SI{0.733}{\gram} of \ce{CO2}
and \SI{0.300}{\gram} of \ce{H2O}. Its molar mass is
\SI{60.05}{\gram\per\mole}.

\begin{parts}
  \item Determine the empirical and molecular formulas.
        \workspace[3.6cm]
        \keyonly{\begin{apcanswer}Carbon, from the \ce{CO2}:
        $n = 0.733/44.01 = \num{0.01665}$~mol C, mass
        $= \SI{0.200}{\gram}$.

        Hydrogen, from the \ce{H2O}: $n(\ce{H2O}) = 0.300/18.02 =
        \num{0.01665}$~mol, so $n(\ce{H}) = \num{0.03331}$~mol, mass
        $= \SI{0.0336}{\gram}$.

        Oxygen \emph{by difference} --- it cannot be taken from the
        products, since the burning oxygen came from the air:
        $0.500 - 0.200 - 0.034 = \SI{0.266}{\gram}$, so
        $n = 0.266/16.00 = \num{0.01665}$~mol.

        Ratio $\ce{C}:\ce{H}:\ce{O} = 1:2:1$, so the empirical formula is
        \textbf{\ce{CH2O}} (mass \num{30.03}).

        $60.05/30.03 = 2$, so the molecular formula is
        \textbf{\ce{C2H4O2}}.\end{apcanswer}}
\end{parts}

\begin{rubric}
  \rpoint{4}{1 pt carbon from \ce{CO2}; 1 pt hydrogen from \ce{H2O}, with
    the factor of 2; 1 pt oxygen \textbf{by difference}; 1 pt both
    formulas. Taking oxygen from the \ce{CO2} and \ce{H2O} masses is the
    classic error and forfeits the third point}
\end{rubric}

% =====================================================================
\frq{short}{4}{\ced{1.4} \textbullet\ composition of a mixture}

A \SI{5.00}{\gram} sample of a sand-and-salt mixture is stirred with
water, filtered, and the filtrate evaporated to dryness, leaving
\SI{3.20}{\gram} of \ce{NaCl}.

\begin{parts}
  \item Calculate the percent by mass of \ce{NaCl} in the mixture.
        \answerlines[2]{$\dfrac{3.20}{5.00} \times 100 =
        \mathbf{64.0}\%$}
  \item Explain why this separation works, in terms of a property that
        differs between the two components.
        \answerlines[3]{The components differ in \textbf{solubility in
        water}. Sodium chloride is ionic and dissolves as hydrated
        \ce{Na+} and \ce{Cl-} ions; sand (\ce{SiO2}) is a covalent network
        solid held by strong covalent bonds throughout, and water cannot
        provide interactions strong enough to break it apart. Filtration
        then separates the undissolved solid from the solution.}
\end{parts}

\begin{rubric}
  \rpoint{2}{(a) 1 pt the ratio; 1 pt \SI{64.0}{\percent}}
  \rpoint{2}{(b) 1 pt naming the solubility difference; 1 pt a reason
    grounded in bonding or particle-level interaction. ``Salt dissolves
    and sand does not'' with no reason earns 1}
\end{rubric}

% =====================================================================
\frq{short}{4}{\ced{1.8} \textbullet\ ionic compounds and Coulombic attraction}

\begin{parts}
  \item Calcium reacts with nitrogen to form a binary ionic compound.
        Give the ions formed and the formula of the compound.
        \answerlines[2]{\ce{Ca^2+} and \ce{N^3-}; the neutral formula is
        \textbf{\ce{Ca3N2}}, since $3(+2) + 2(-3) = 0$.}
  \item Predict which has the larger lattice energy, \ce{CaO} or
        \ce{KCl}, and justify in Coulombic terms. \workspace[2.6cm]
        \keyonly{\begin{apcanswer}\textbf{\ce{CaO}}. The Coulombic
        attraction is proportional to the product of the ionic charges:
        $(+2)(-2) = 4$ for \ce{CaO} against $(+1)(-1) = 1$ for \ce{KCl}.

        \ce{Ca^2+} and \ce{O^2-} are also the smaller ions, so the
        internuclear distance is shorter and the attraction, which varies
        as $1/d$, is stronger again. Both factors point the same
        way.\end{apcanswer}}
\end{parts}

\begin{rubric}
  \rpoint{2}{(a) 1 pt both ions; 1 pt \ce{Ca3N2}}
  \rpoint{2}{(b) 1 pt \ce{CaO}; 1 pt a Coulombic justification naming
    charge (and ideally distance). ``\ce{CaO} is more stable'' or an
    octet-based answer earns \textbf{0}}
\end{rubric}

\examtotal

\end{document}
